\section{Review of Related Systems}
This section explores prominent interlocking systems that have shaped the industry and influenced our design.

\subsection{European Train Control System (ETCS)}
ETCS is a standardized signaling system under the European Rail Traffic Management System (ERTMS). It integrates interlocking with cab signaling and automatic train protection. ETCS Level 2 and 3 eliminate reliance on trackside signals, relying on continuous communication with onboard equipment. While highly advanced, it is cost-intensive and primarily suited for high-speed corridors.

\subsection{Solid State Interlocking (SSI)}
SSI was one of the first commercial microprocessor-based interlocking systems. Deployed in the UK and Europe, it uses a centralized computer to process logic and interface with relay interfaces. It demonstrated early benefits of programmability and space efficiency but was limited in graphical diagnostics and simulation support.

\subsection{Bangladesh Railway Practices}
Bangladesh Railways still predominantly relies on token-based or relay-based signaling, especially in non-electrified regions. While some modernization efforts have introduced computerized panels, the interlocking logic remains largely static and manually supervised. This environment highlights the need for cost-effective, software-centric solutions that are deployable in constrained settings.
